\section{Covariance Dynamics and Equilibrium Analysis}
\label{app:covariance_dynamics}

\paragraph{Scope of this appendix.}
The analysis below is a local, fixed-attention, receiver-only surrogate. We treat the attention weights $\beta_{ij}$ as exogenous constants when differentiating with respect to $\Sigma_i$, vary only the receiver covariance, and never differentiate the sender covariances $\Sigma_j$ or the weights $\beta_{ij}$ themselves. The full reduced functional, obtained by envelope-eliminating $\beta_{ij}$ via its entropy-regularized softmax optimum and including the complete sender-row contributions, is not derived here; no explicit first-derivative $\partial\beta_{ij}$ response survives in that canonical functional. The fixed point characterized below is therefore a partial Euler-Lagrange stationarity condition, not a stationary point of the fully reduced free energy.

The covariance fixed point is also a property of the consensus potential of Section~\ref{sec:mixture_derivation}, which weights the belief entropy of agent $i$ by $1+B_i$ with $B_i := \sum_j\beta_{ij}$. Minimizing $\mathbb{E}_{q}[U] - (1+B_i)H(q)$ over laws $q$ gives $q \propto e^{-U/(1+B_i)}$, so each belief is a Gibbs state at temperature $1+B_i = 2$ in units in which Bayesian updating has temperature one. In a normalized Gaussian joint the entropy enters once, and the mean-field covariance is the inverse diagonal precision block, fixed in one step and independent of the neighbors' covariances~\cite[Ch.~G4, block optimum]{DennisGaugeVFESuppG2026}. The factor $\tfrac{1}{2}$ in Eq.~\eqref{eq:sigma_fixed_point_beta} and the homogeneous limit $\Sigma_\infty = \Sigma_0$ below both come from this doubled entropy weight.

\subsection{Covariance Gradient of the Per-Agent Block}

Here we derive the gradient of agent $i$'s block $\mathcal{F}_i$ of the consensus functional with respect to its covariance $\Sigma_i$, a standard computation for Gaussian families~\cite{Amari2016, petersen2012matrix}.

The block is

\begin{equation}
\mathcal{F}_i
=D_{\mathrm{KL}}(q_i \| p_i)
+ \sum_{j \neq i} \beta_{ij}
D_{\mathrm{KL}}(q_i \| \Omega_{ij} q_j)
- \mathbb{E}_{q_i}[\log p(o_i \mid k_i)],
\end{equation}

where $q_i = \mathcal{N}(\mu_i, \Sigma_i)$, $p_i = \mathcal{N}(\mu_{p,i}, \Sigma_{p,i})$, and $\sum_j \beta_{ij} = 1$ by construction.

\subsubsection{Gaussian relative entropy and its derivative}

Differentiating the Gaussian relative entropy of Appendix~\ref{app:gaussian_kl}, with first argument $\mathcal{N}(\mu_1,\Sigma_1)$ and second argument $\mathcal{N}(\mu_2,\Sigma_2)$, with respect to $\Sigma_1$ at fixed $\mu_1$, $\mu_2$, $\Sigma_2$ gives

\begin{equation}
\frac{\partial D_{\mathrm{KL}}}{\partial \Sigma_1}
=\frac{1}{2}\left[
-\Sigma_1^{-1}
+ \Sigma_2^{-1}
\right].
\end{equation}

Applying this to each term in $\mathcal{F}_i$:

\begin{align}
\frac{\partial}{\partial \Sigma_i}
D_{\mathrm{KL}}(q_i \| p_i)
&=\frac{1}{2}\left[
-\Sigma_i^{-1}
+ \Sigma_{p,i}^{-1}
\right],
\\[6pt]
\frac{\partial}{\partial \Sigma_i}
D_{\mathrm{KL}}(q_i \| \Omega_{ij} q_j)
&=
\frac{1}{2}\left[
-\Sigma_i^{-1}
+ (\Omega_{ij}\Sigma_j\Omega_{ij}^\top)^{-1}
\right].
\end{align}

The observation term, for the Gaussian likelihood $-\mathbb{E}_{q_i}[\log p(o_i \mid k_i)] = \tfrac{1}{2}(o_i - \mu_i)^\top \Lambda_{o_i}(o_i - \mu_i) + \tfrac{1}{2}\mathrm{tr}(\Lambda_{o_i}\Sigma_i) + \mathrm{const}$, contributes a constant precision correction $\partial[-\mathbb{E}_{q_i}\log p(o_i \mid k_i)]/\partial\Sigma_i = \tfrac{1}{2}\Lambda_{o_i}$, independent of $\Sigma_i$. Including it,

\begin{equation}
\boxed{
\frac{\partial \mathcal{F}_i}{\partial \Sigma_i}
=\frac{1}{2}
\left[
-2 \Sigma_i^{-1}
+ \Sigma_{p,i}^{-1}
+ \sum_j \beta_{ij}
(\Omega_{ij}\Sigma_j\Omega_{ij}^\top)^{-1}
+ \Lambda_{o_i}
\right].
}
\end{equation}

Because $\sum_j \beta_{ij} = 1$, the prefactor $(1 + \sum_j \beta_{ij})$ in front of $\Sigma_i^{-1}$ reduces to the constant $2$, as displayed in the boxed gradient.

\subsection{Fixed-Point Equation and Symmetric Solution}

At equilibrium, we set $\partial \mathcal{F}_i / \partial \Sigma_i = 0$, giving

\begin{equation}
\Sigma_i^{-1}
=\frac{1}{2}
\left[
\Sigma_{p,i}^{-1}
+ \sum_j \beta_{ij}
(\Omega_{ij}\Sigma_j\Omega_{ij}^\top)^{-1}
+ \Lambda_{o_i}
\right].
\label{eq:sigma_fixed_point_beta}
\end{equation}
The sensory precision $\Lambda_{o_i}$ enters with the same prefactor as the other precision contributions; setting $\Lambda_{o_i} = 0$ gives the observation-free case used in the homogeneous limit below. Equation~\eqref{eq:sigma_fixed_point_beta} couples all agents: each agent's precision is half the sum of its prior and sensory precision and the $\beta$-weighted transported neighbor precisions.

\paragraph{Precision coordinates.} Write $\Lambda_i := \Sigma_i^{-1}$ for the belief precision, $A_i := \Sigma_{p,i}^{-1} + \Lambda_{o_i} \succ 0$ for the prior-and-sensory anchor, and define the linear map
\begin{equation}
\mathcal{L}(\Lambda)_i := \sum_j \beta_{ij}\Omega_{ij}^{-\top}\Lambda_j\Omega_{ij}^{-1}
\label{eq:precision_transport_map}
\end{equation}
on the space $\mathrm{Sym}(K)^N$ of $N$-tuples of symmetric matrices. Since $(\Omega_{ij}\Sigma_j\Omega_{ij}^\top)^{-1} = \Omega_{ij}^{-\top}\Lambda_j\Omega_{ij}^{-1}$, Eq.~\eqref{eq:sigma_fixed_point_beta} is linear in precision coordinates,
\begin{equation}
2\Lambda = A + \mathcal{L}(\Lambda),
\label{eq:precision_linear}
\end{equation}
and the covariance fixed-point map
\begin{equation}
T(\Sigma)_i = 2\big(\Sigma_{p,i}^{-1} + \sum_j \beta_{ij}(\Omega_{ij}\Sigma_j\Omega_{ij}^\top)^{-1} + \Lambda_{o_i}\big)^{-1}
\label{eq:sigma_fixed_point_map}
\end{equation}
is, in precision coordinates, the affine map $\Lambda \mapsto \tfrac{1}{2}(A + \mathcal{L}(\Lambda))$. Each term of $\mathcal{L}$ is a congruence with a nonnegative weight, so $\mathcal{L}$ carries the product cone of positive semidefinite tuples into itself. Under a change of frames, $\Lambda_i \mapsto g_i^{-\top}\Lambda_i g_i^{-1}$ and $\Omega_{ij} \mapsto g_i\Omega_{ij}g_j^{-1}$, the map $\mathcal{L}$ transforms by similarity, so its spectral radius $\rho(\mathcal{L})$ is frame independent.

\begin{proposition}[Frozen-attention covariance fixed point]
\label{prop:covariance_fixed_point}
Fix $\beta_{ij} \geq 0$ and $\Omega_{ij} \in \mathrm{GL}(K)$, and assume $\Sigma_{p,i} \succ 0$ and $\Lambda_{o_i} \succeq 0$.
\begin{enumerate}
\item[(i)] Equation~\eqref{eq:sigma_fixed_point_beta} has a solution with every $\Sigma_i \succ 0$ if and only if $\rho(\mathcal{L}) < 2$. The solution is then unique, $\Lambda^* = \tfrac{1}{2}\sum_{k \geq 0}(\mathcal{L}/2)^k A$, and the iteration of $T$ converges to it geometrically with asymptotic ratio $\rho(\mathcal{L})/2$.
\item[(ii)] If $\rho(\mathcal{L}) < 2$, the simultaneous receiver-only natural-gradient flow of all agents, $d\Sigma_i/dt = -2\eta_\Sigma\Sigma_i(\partial\mathcal{F}_i/\partial\Sigma_i)\Sigma_i$ with $\eta_\Sigma > 0$, remains in the SPD cone and converges to $\Lambda^*$ from every SPD initial condition, at the exponential rate $2\eta_\Sigma(1 - \rho(\mathcal{L})/2)$.
\item[(iii)] If $\rho(\mathcal{L}) \geq 2$, no SPD fixed point exists, and along the same flow the precision grows without bound, so that some covariance direction collapses.
\item[(iv)] $\rho(\mathcal{L}) \leq \max_i\lambda_{\max}(C_i)$ with $C_i := \sum_j\beta_{ij}(\Omega_{ij}\Omega_{ij}^\top)^{-1} = \mathcal{L}(I)_i$. If every row sums to one and a change of frames makes every link orthogonal, in particular if they are graph-flat, $\Omega_{ij} = U_iU_j^{-1}$, then $\rho(\mathcal{L}) = 1$.
\end{enumerate}
\end{proposition}

\begin{proof}
With the boxed gradient, $\partial\mathcal{F}_i/\partial\Sigma_i = \tfrac{1}{2}(A_i + \mathcal{L}(\Lambda)_i - 2\Lambda_i)$, and $d\Lambda_i/dt = -\Lambda_i(d\Sigma_i/dt)\Lambda_i$, so the flow of (ii) is the linear system
\begin{equation}
\frac{d\Lambda}{dt} = \eta_\Sigma\big(A + \mathcal{L}(\Lambda) - 2\Lambda\big).
\label{eq:precision_flow}
\end{equation}
For (i), if $\rho(\mathcal{L}) < 2$ the Neumann series converges, $2I - \mathcal{L}$ is invertible, and $\Lambda^* \succeq A/2 \succ 0$ because every term of the series is positive semidefinite; the iteration error is $(\mathcal{L}/2)^n(\Lambda_0 - \Lambda^*)$. Conversely, let $\Lambda \succ 0$ solve Eq.~\eqref{eq:precision_linear}, and choose $\epsilon > 0$ with $A_i \succeq \epsilon\Lambda_i$ for every $i$. Then $\mathcal{L}(\Lambda) = 2\Lambda - A \preceq (2-\epsilon)\Lambda$. In the order-unit norm $\|X\|_\Lambda := \min\{s \geq 0 : -s\Lambda \preceq X \preceq s\Lambda\}$, which is a norm because $\Lambda$ is interior to the cone, cone preservation gives $\|\mathcal{L}(X)\|_\Lambda \leq (2-\epsilon)\|X\|_\Lambda$, hence $\rho(\mathcal{L}) \leq 2 - \epsilon < 2$. The same argument with $\Lambda = I$ and $\mathcal{L}(I) \preceq \max_i\lambda_{\max}(C_i) I$ proves the bound in (iv). For (ii), variation of constants gives
\[
\Lambda(t) = e^{-2\eta_\Sigma t}e^{\eta_\Sigma t\mathcal{L}}\Lambda(0) + \eta_\Sigma\int_0^t e^{-2\eta_\Sigma s}e^{\eta_\Sigma s\mathcal{L}}(A)\mathrm{d}s.
\]
The exponential $e^{s\mathcal{L}} = \sum_k (s\mathcal{L})^k/k!$ preserves the cone and satisfies $e^{s\mathcal{L}}A \succeq A$, so $\Lambda(t) \succeq (1 - e^{-2\eta_\Sigma t})A/2$ plus a positive semidefinite term, and $\Lambda(t)$ stays SPD. The deviation is $\Lambda(t) - \Lambda^* = e^{-\eta_\Sigma(2I - \mathcal{L})t}(\Lambda(0) - \Lambda^*)$. Every eigenvalue of $2I - \mathcal{L}$ has real part at least $2 - \rho(\mathcal{L})$, with equality at $2 - \rho(\mathcal{L})$ itself, because the spectral radius of a cone-preserving map is an eigenvalue~\cite{Vandergraft1968}. For (iii), the same theorem applied to the adjoint $\mathcal{L}^*(W)_j = \sum_i\beta_{ij}\Omega_{ij}^{-1}W_i\Omega_{ij}^{-\top}$, which preserves the positive semidefinite cone (self-dual under $\langle X, Y\rangle = \sum_i\mathrm{tr}(X_iY_i)$), supplies $W \succeq 0$, $W \neq 0$, with $\mathcal{L}^*W = \rho(\mathcal{L})W$. Pairing a fixed point with $W$ gives $(2 - \rho(\mathcal{L}))\langle W, \Lambda\rangle = \langle W, A\rangle > 0$, which is impossible for $\rho(\mathcal{L}) \geq 2$ and $\Lambda \succ 0$. Along Eq.~\eqref{eq:precision_flow}, $d\langle W, \Lambda\rangle/dt = \eta_\Sigma(\langle W, A\rangle + (\rho(\mathcal{L}) - 2)\langle W, \Lambda\rangle) \geq \eta_\Sigma\langle W, A\rangle > 0$, since $\Lambda(t)$ stays positive semidefinite, so $\langle W, \Lambda(t)\rangle$ grows at least linearly in $t$. For the second part of (iv), if frames $g_i$ make every link orthogonal, then in those frames $\mathcal{L}(I)_i = \sum_j\beta_{ij}I = I$, so $I$ is an interior eigenvector with eigenvalue one and the bound gives $\rho(\mathcal{L}) = 1$; similarity carries this back to the original frames. For graph-flat links the frames $g_i = U_i^{-1}$ make every link the identity.
\end{proof}

\paragraph{Frame dependence and holonomy.} A Thompson-metric contraction argument for the map $T$~\cite{nussbaum1988hilbert, lawson2001geometric} yields the same fixed point under the hypothesis $\lambda_{\max}(C_i) < 2$ for every $i$. By part (iv) that hypothesis is sufficient, but it is not frame independent, because $C_i$ picks up the factor $g_j^{-1}g_j^{-\top}$ under a change of frames. For graph-flat $\mathrm{GL}(K)$ links it can fail by a large margin while $\rho(\mathcal{L}) = 1$: for three agents with frames $U = \mathrm{diag}(10, 1)$, $I$, and $\bigl(\begin{smallmatrix}1 & 3\\ 0 & 1/2\end{smallmatrix}\bigr)$, one finds $\max_i\lambda_{\max}(C_i) \approx 610$, while the fixed point exists and is globally stable. The value of $\rho(\mathcal{L})$ departs from one only through loop holonomy whose holonomy group is not conjugate into $\mathrm{O}(K)$, and such holonomy can push it past two. For $K = 1$ and two agents with $\beta_{12} = \beta_{21} = 1$ and independent links $w = \Omega_{12}$, $v = \Omega_{21}$, the map is $\mathcal{L} = \bigl(\begin{smallmatrix}0 & w^{-2}\\ v^{-2} & 0\end{smallmatrix}\bigr)$ with $\rho(\mathcal{L}) = 1/h$, where $h = wv$ is the backtrack holonomy; no SPD fixed point exists for $h \leq 1/2$, and the iterated covariances decay to zero. With orientation reversal $\Omega_{ji} = \Omega_{ij}^{-1}$ imposed, as for the independent links of Regime~II, a two-cycle carries no holonomy, but the obstruction reappears on a triangle: with rows $\beta_{ij} = 1/2$ and holonomy $h$ around the loop, the characteristic polynomial of $\mathcal{L}$ is $\lambda^3 - \tfrac{3}{4}\lambda - (h^2 + h^{-2})/8$, and $\rho(\mathcal{L}) \geq 2$ exactly when $h^2 + h^{-2} \geq 52$.

\subsubsection{Homogeneous limit}

In the homogeneous limit we assume that (i) all agents are identical, so $\Sigma_i = \Sigma_\infty$ for all $i$; (ii) every link preserves the shared prior covariance, $\Omega_{ij}\Sigma_0\Omega_{ij}^\top = \Sigma_0$ (for instance $\Omega_{ij} = I$); (iii) $\Sigma_{p,i} = \Sigma_0$; and (iv) $\Lambda_{o_i} = 0$, the observation-free case. Under (i)-(iv), \eqref{eq:sigma_fixed_point_beta} becomes

\begin{equation}
\Sigma_\infty^{-1}
=\frac{1}{2}
\left[
\Sigma_0^{-1}
+ \Sigma_\infty^{-1}
\right]
\quad\Longrightarrow\quad
\Sigma_\infty = \Sigma_0,
\end{equation}
and the solution is unique by Proposition~\ref{prop:covariance_fixed_point}, since links preserving $\Sigma_0$ are orthogonal in the common frame $\Sigma_0^{-1/2}$. With a common sensory precision $\Lambda_{o_i} = \Lambda_o > 0$ and links that also preserve $\Sigma_0^{-1} + \Lambda_o$, the same computation gives $\Sigma_\infty = (\Sigma_0^{-1} + \Lambda_o)^{-1}$.

Hence, in a symmetric observation-free population, the equilibrium covariance reproduces the shared prior. This reflects the doubled entropy weight of the consensus potential: a normalized Gaussian joint with prior precision $\Sigma_0^{-1}$ and edge precision $\Sigma_0^{-1}$ gives the mean-field covariance $\Sigma_0/2$, which is precision addition~\cite[Ch.~G4, block optimum]{DennisGaugeVFESuppG2026}.

\subsubsection{Alignment-dominated regime}

Although $\sum_j \beta_{ij} = 1$, one might expect the effective strength of alignment to be controlled by the attention temperature $\tau$ in

\begin{equation}
\beta_{ij}
=\frac{
\exp\left[
-\frac{1}{\tau}
D_{\mathrm{KL}}\big(q_i \| \Omega_{ij} q_j\big)
\right]
}{\sum_k
\exp\left[
-\frac{1}{\tau}
D_{\mathrm{KL}}\big(q_i \| \Omega_{ik} q_k\big)
\right]}.
\end{equation}

As $\tau \to 0$, $\beta_{ij}$ becomes sharply peaked on $j^\star = \arg\min_j D_{\mathrm{KL}}(q_i \| \Omega_{ij}q_j)$. Since $\sum_j \beta_{ij} = 1$ for every $\tau > 0$, sharpening $\beta$ by itself does not amplify the coefficient on the alignment block: in the limit, $\sum_j \beta_{ij}(\Omega_{ij}\Sigma_j\Omega_{ij}^\top)^{-1} \to (\Omega_{ij^\star}\Sigma_{j^\star}\Omega_{ij^\star}^\top)^{-1}$, which is bounded. To make the alignment-dominated regime precise we place an explicit coupling strength $a > 0$ on the alignment block,

\begin{equation}
\mathcal{F}_i^{(a)}
=D_{\mathrm{KL}}(q_i \| p_i)
+ a \sum_{j \neq i} \beta_{ij}
D_{\mathrm{KL}}(q_i \| \Omega_{ij} q_j)
- \mathbb{E}_{q_i}[\log p(o_i \mid k_i)],
\end{equation}

of which the derivation above is the case $a = 1$. Repeating the differentiation term by term, the stationarity condition becomes

\begin{equation}
(1 + a)\Sigma_i^{-1}
=\Sigma_{p,i}^{-1}
+ a \sum_j \beta_{ij}
(\Omega_{ij}\Sigma_j\Omega_{ij}^\top)^{-1}
+ \Lambda_{o_i},
\end{equation}

which reduces to \eqref{eq:sigma_fixed_point_beta} at $a = 1$. In precision coordinates it reads $(1+a)\Lambda = A + a\mathcal{L}(\Lambda)$, and the argument of Proposition~\ref{prop:covariance_fixed_point}(i), applied to $a\mathcal{L}/(1+a)$, shows that an SPD solution exists if and only if $\rho(\mathcal{L}) < 1 + 1/a$. The high-coupling limit $a \to \infty$ therefore exists only if $\rho(\mathcal{L}) \leq 1$. Dividing by $a$ gives formally

\begin{equation}
\boxed{
\Sigma_i^{-1}
\longrightarrow
\sum_j \beta_{ij}
(\Omega_{ij}\Sigma_j\Omega_{ij}^\top)^{-1}
=\big\langle(\Omega_{ij}\Sigma_j\Omega_{ij}^\top)^{-1}
\big\rangle_{\beta}
\qquad \text{(alignment-dominated)},}
\label{eq:beta_weighted_precision}
\end{equation}

where $\langle \cdot \rangle_{\beta}$ denotes a $\beta$-weighted average over neighbors. This is the eigenvalue equation $\Lambda = \mathcal{L}(\Lambda)$. If $\rho(\mathcal{L}) < 1$, the finite-coupling solution is $\Lambda_a = (I + a(I - \mathcal{L}))^{-1}A$, which tends to zero like $1/a$, so the covariances diverge. The boxed relation has a nonzero SPD limit only when $\rho(\mathcal{L}) = 1$ with an SPD Perron eigenvector. Graph-flat links are the basic case. In the frames $g_i = U_i^{-1}$ the map acts as $\tilde\Lambda_i \mapsto \sum_j\beta_{ij}\tilde\Lambda_j$ with $\tilde\Lambda_i := U_i^\top\Lambda_iU_i$. For an irreducible row-stochastic $\beta$ with stationary distribution $\pi$, contracting the finite-coupling equation with $\pi$ gives $\sum_i\pi_i\tilde\Lambda_i = \sum_i\pi_i\tilde A_i$ for every $a$, with $\tilde A_i := U_i^\top A_iU_i$, while the components off the Perron direction decay like $1/a$. The limit is therefore the common precision $\tilde\Lambda_\infty = \sum_k\pi_k\tilde A_k$, and the limit covariances $\Sigma_i = U_i\tilde\Lambda_\infty^{-1}U_i^\top$ satisfy $\Sigma_i = \Omega_{ij}\Sigma_j\Omega_{ij}^\top$ on every edge.

Per-edge covariance alignment is thus exact in the flat-link, high-coupling limit, and small $\tau$ alone does not produce it. At finite coupling, or with holonomy, the equilibrium satisfies only the $\beta$-averaged relation~\eqref{eq:sigma_fixed_point_beta}, and with several heterogeneous neighbors the per-edge mismatch is generically nonzero, since $\Sigma_i$ cannot equal every transported neighbor covariance at once.

\subsubsection{Stability}

Proposition~\ref{prop:covariance_fixed_point}(ii) gives global stability of the simultaneous receiver-only natural-gradient flow whenever $\rho(\mathcal{L}) < 2$. The own-block Hessian is also exact,
\begin{equation}
\frac{\partial^2 \mathcal{F}_i}{\partial \Sigma_i \partial \Sigma_i}[X]
= \Sigma_i^{-1} X \Sigma_i^{-1},
\end{equation}
since each KL term contributes $\tfrac{1}{2}\Sigma_i^{-1}X\Sigma_i^{-1}$ from its log-determinant, the self and alignment weights sum to $2 \times \tfrac{1}{2} = 1$, and the trace and observation terms are linear in $\Sigma_i$. It is positive definite for every $\Sigma_i \succ 0$, so each agent's problem with its neighbors frozen is strictly convex in $\Sigma_i$. Neither statement covers the complete canonical objective, whose Hessian carries sender cross-blocks and the reduced-Hessian softmax covariance correction; stability of that coupled equilibrium, and existence of an equilibrium of the live-softmax dynamics, are not established here. A Euclidean flow $d\Sigma_i/dt = -\eta_\Sigma\partial\mathcal{F}_i/\partial\Sigma_i$ has the same fixed point, but its linearization depends on the conditioning of the fixed point, and we have no general stability theorem for it.

\paragraph{Mean update where covariances align.}
The natural-gradient mean update $\Sigma_i\nabla_{\mu_i}\mathcal{F}_i$ carries the per-edge alignment term $\beta_{ij} \Sigma_i (\Omega_{ij}\Sigma_j\Omega_{ij}^\top)^{-1}(\mu_i - \Omega_{ij}\mu_j)$. Where per-edge alignment $\Sigma_i = \Omega_{ij}\Sigma_j\Omega_{ij}^\top$ holds, as in the flat-link, high-coupling limit above, the precision factors cancel and
\begin{equation}
\beta_{ij} \Sigma_i (\Omega_{ij}\Sigma_j\Omega_{ij}^\top)^{-1}(\mu_i - \Omega_{ij}\mu_j) \longrightarrow \beta_{ij}(\mu_i - \Omega_{ij}\mu_j),
\end{equation}
so the step is an attention-weighted sum of transported residuals with no covariance factor. This is again a property of the consensus potential: in the exact receiver update of a normalized comparison-record model the corresponding step keeps a Kalman-type gain set by the receiver's posterior covariance~\cite[Ch.~6, exact receiver update]{DennisGaugeVFE2026}.

\section{Relating the Quadratic Forms to Transported Relative Entropies}
We now show that the pairwise quadratic expectations can be written as transported relative entropies. The direction of the divergence follows from building the coupling precision from the transported sender covariance, as shown below; it is not supplied by the representation theorem of Appendix~\ref{app:conditional_representation}.
\subsection{Exact expansion for Gaussian beliefs}
For independent Gaussians $q_i = \mathcal{N}(\mu_{q,i},\Sigma_{q,i})$ and $q_j = \mathcal{N}(\mu_{q,j},\Sigma_{q,j})$ and any symmetric matrix $M$, the expectation of a quadratic form is
\begin{equation}
\mathbb{E}_{q_i q_j}[\delta^\top M \delta]
= \mathrm{tr}\big(M\mathrm{Cov}(\delta)\big)
+ \bar{\delta}^\top M \bar{\delta},
\label{eq:quadratic_expectation_formula}
\end{equation}
where $\delta = k_i - \Omega_{ij}k_j$, $\bar{\delta} = \mathbb{E}[\delta] = \mu_{q,i} - \Omega_{ij}\mu_{q,j}$, and
\begin{equation}
\mathrm{Cov}(\delta) = \Sigma_{q,i} + \Omega_{ij}\Sigma_{q,j}\Omega_{ij}^\top.
\end{equation}
Applying \eqref{eq:quadratic_expectation_formula} with $M = W^{(q)}_{ij}$:
\begin{equation}
\begin{aligned}
\mathbb{E}_{q_i q_j}
[(k_i-\Omega_{ij}k_j)^\top W^{(q)}_{ij}(k_i-\Omega_{ij}k_j)]
&=\mathrm{tr}\big[W^{(q)}_{ij}(\Sigma_{q,i}+\Omega_{ij}\Sigma_{q,j}\Omega_{ij}^\top)\big]\\
&\quad+(\mu_{q,i}-\Omega_{ij}\mu_{q,j})^\top
W^{(q)}_{ij}
(\mu_{q,i}-\Omega_{ij}\mu_{q,j}).
\end{aligned}
\label{eq:belief_quadratic_expansion}
\end{equation}
An identical expansion holds for the model channel with $W^{(s)}_{ij}$ and $\tilde{\Omega}_{ij}$. For a fixed coupling $W_{ij} \succ 0$ the mean-dependent part is exactly a relative entropy at common covariance: $\tfrac{1}{2}\mathbb{E}[\delta^\top W_{ij}\delta] = D_{\mathrm{KL}}\big(\mathcal{N}(\mu_{q,i}, W_{ij}^{-1})\|\mathcal{N}(\Omega_{ij}\mu_{q,j}, W_{ij}^{-1})\big) + \tfrac{1}{2}\mathrm{tr}\big(W_{ij}(\Sigma_{q,i} + \Omega_{ij}\Sigma_{q,j}\Omega_{ij}^\top)\big)$~\cite[Ch.~G5, edge relative entropy at fixed covariance]{DennisGaugeVFESuppG2026}.

To obtain the transported relative entropy of the beliefs themselves we choose state-dependent couplings,
\begin{equation}
W^{(q)}_{ij}:=c^{(q)}_{ij}(\Omega_{ij}\Sigma_{q,j}\Omega_{ij}^\top)^{-1},
\qquad
W^{(s)}_{ij}:=c^{(s)}_{ij}(\tilde{\Omega}_{ij}\Sigma_{s,j}\tilde{\Omega}_{ij}^\top)^{-1},
\end{equation}
where $c^{(q)}_{ij}, c^{(s)}_{ij} > 0$ are dimensionless coupling strengths. Because $W^{(q)}_{ij}$ is built from the current variational covariance $\Sigma_{q,j}$ (and $W^{(s)}_{ij}$ from $\Sigma_{s,j}$), the identification below is a value-level matching at the current variational state, not a property of a fixed generative model, consistent with the status of the construction in Section~\ref{sec:mixture_derivation}. No dynamics in this manuscript are obtained by differentiating the quadratic form with $W^{(q)}_{ij}$ held fixed; all gradient flows differentiate the KL coupling directly.

Inserting the Gaussian KL closed form of Appendix~\ref{app:gaussian_kl} for $q_i = \mathcal{N}(\mu_{q,i},\Sigma_{q,i})$ and $\Omega_{ij}q_j = \mathcal{N}(\Omega_{ij}\mu_{q,j}, \Omega_{ij}\Sigma_{q,j}\Omega_{ij}^\top)$ into \eqref{eq:belief_quadratic_expansion} gives the exact identity
\begin{equation}
\begin{aligned}
\frac{1}{4}\mathbb{E}_{q_i q_j}
[(k_i-\Omega_{ij}k_j)^\top W^{(q)}_{ij}(k_i-\Omega_{ij}k_j)]
&=
\frac{c^{(q)}_{ij}}{2}D_{\mathrm{KL}}(q_i\|\Omega_{ij}q_j)+\frac{c^{(q)}_{ij}K_q}{2}\\
&\quad+\frac{c^{(q)}_{ij}}{4}\log\det\big((\Omega_{ij}\Sigma_{q,j}\Omega_{ij}^\top)^{-1}\Sigma_{q,i}\big).
\end{aligned}
\label{eq:quadratic_kl_identity}
\end{equation}
The identification with the transported relative entropy is thus exact up to the constant and the single log-determinant term. The direction comes from the transported sender covariance in $W^{(q)}_{ij}$: with $W = c\Sigma_{q,i}^{-1}$ the same expansion yields $D_{\mathrm{KL}}(\Omega_{ij}q_j\|q_i)$. The log-determinant term vanishes when $\det\Sigma_{q,i} = \det(\Omega_{ij}\Sigma_{q,j}\Omega_{ij}^\top)$. Expanded in the covariance mismatch $\Delta := \Sigma_{q,i} - \Omega_{ij}\Sigma_{q,j}\Omega_{ij}^\top$, its leading contribution is $\tfrac{c^{(q)}_{ij}}{4}\mathrm{tr}\big((\Omega_{ij}\Sigma_{q,j}\Omega_{ij}^\top)^{-1}\Delta\big)$, so the correction is $O(\|\Delta\|)$ in general and $O(\|\Delta\|^2)$ only along determinant-preserving deviations. With several heterogeneous neighbors the per-edge mismatch $\Delta_{ij}$ is generically nonzero at the stationary point of Eq.~\eqref{eq:sigma_fixed_point_beta} (Appendix~\ref{app:covariance_dynamics}), and the first-order correction persists there. The model channel satisfies the same identity with $(m, s, \tilde\Omega_{ij}, W^{(s)}_{ij}, c^{(s)}_{ij}, K_m)$ in place of $(k, q, \Omega_{ij}, W^{(q)}_{ij}, c^{(q)}_{ij}, K_q)$.
\subsection{Defining raw alignment weights}
\label{sec:raw_alignment_weights}
To connect with the main text we define the raw alignment weights
\begin{equation}
\tilde\beta_{ij}:=\frac{c^{(q)}_{ij}}{2},
\qquad
\tilde\gamma_{ij}:=\frac{c^{(s)}_{ij}}{2},
\end{equation}
the bare belief- and model-alignment strengths between agents $i$ and $j$. In the limit $\tilde\beta_{ij} \to \infty$ with fixed $\tilde\gamma_{ij}$, agents' beliefs are forced to agree after transport while their models may still differ; $\tilde\gamma_{ij} = 0$ decouples model alignment entirely. The normalized weights $\beta_{ij}, \gamma_{ij}$ of the main text, which sum to one over each row, are not identified with the raw couplings. They arise by composition: at a uniform coupling $c^{(q)}_{ij} \equiv c^{(q)}$, the identity~\eqref{eq:quadratic_kl_identity} supplies the per-edge energies $E_{ij} = D_{\mathrm{KL}}(q_i \| \Omega_{ij}q_j)$ up to the common prefactor, and minimizing the entropy-regularized row objective of Section~\ref{sec:mixture_derivation} (Eqs.~\eqref{eq:softmax_attention_general} and~\eqref{eq:beta_optimal}) over the row simplex returns the softmax weights. The quadratic form reproduces the canonical functional only through this composition, and only up to the log-determinant term of Eq.~\eqref{eq:quadratic_kl_identity}.

\section{Conditional Representation Theorem for the Transported Relative Entropy}
\label{app:conditional_representation}

This appendix characterizes the belief coupling $D_{\mathrm{KL}}(q_i\|\Omega_{ij}q_j)$, with the receiver law in the first slot (the variational or I-projection direction), within a class of divergences and for one problem: agent $i$'s best response with its neighbors' beliefs and its attention weights held fixed. The companion theory declares this transported divergence and does not derive it~\cite[Chs.~6 and~12]{DennisGaugeVFE2026}, and nothing below fixes the divergence of the joint scalar functional; the scope paragraph at the end of this appendix shows why the closure criterion cannot. Within the class of local $f$-divergences that enter the free energy linearly in $\beta_{ij}$, the relative entropy is the only divergence whose best response has the geometric-mean Boltzmann form for all priors and neighbor beliefs; richness of the attainable likelihood ratios, which the converse needs, holds automatically in the Gaussian fiber (Lemma~\ref{lem:richness_app}). We assume matched fibers, a common dimension $K_q$ and representation $\rho$; the argument uses only the density of the transported law $\Omega_{ij}q_j$ and is unchanged when an intertwiner produces it.

\subsection*{The best-response problem}

With the neighbor beliefs $q_j$ and the weights $\beta_{ij} \geq 0$ held fixed, agent $i$ solves
\begin{equation}
F_i[\beta_i]
=\min_{q_i}
\Big\{
D_{\mathrm{KL}}(q_i \| p_i)
+ \sum_{j\neq i} \beta_{ij}\mathcal{D}(q_i,\Omega_{ij}q_j)
\Big\},
\qquad
\mathcal{D}(q, r) = \int r(c) f\left(\frac{q(c)}{r(c)}\right)dc,
\label{eq:F_i_app}
\end{equation}
with the second-argument density as the integration measure. The weights are in turn chosen by minimizing $\sum_{j}\beta_{ij}C_{ij} + \tau\sum_{j}\beta_{ij}\log\beta_{ij}$ over the row simplex, where $C_{ij} := \partial F_i/\partial\beta_{ij}$ is the marginal cost of attending to agent $j$.

\begin{lemma}[KL-sum identity]
\label{lem:kl_sum_identity}
Let $p$ and $r_1, \dots, r_n$ be probability densities with overlapping supports, let $\beta_j \geq 0$ with $B := \sum_j\beta_j$, and define
\begin{equation}
g := Z_g^{-1}p^{1/(1+B)}\prod_j r_j^{\beta_j/(1+B)},
\qquad
Z_g := \int p^{1/(1+B)}\prod_j r_j^{\beta_j/(1+B)}dc.
\end{equation}
Then for every density $q$ with $D_{\mathrm{KL}}(q\|p)$ and every $D_{\mathrm{KL}}(q\|r_j)$ finite,
\begin{equation}
D_{\mathrm{KL}}(q\|p) + \sum_j\beta_jD_{\mathrm{KL}}(q\|r_j) = (1+B)D_{\mathrm{KL}}(q\|g) - (1+B)\log Z_g.
\label{eq:kl_sum_identity}
\end{equation}
\end{lemma}

\begin{proof}
The left-hand side is $\int q[(1+B)\log q - \log p - \sum_j\beta_j\log r_j]dc = (1+B)\int q[\log q - \log g - \log Z_g]dc$, which is the right-hand side.
\end{proof}

With $r_j = \Omega_{ij}q_j$ and $\mathcal{D} = D_{\mathrm{KL}}$, the identity solves the best-response problem~\eqref{eq:F_i_app} at once: the unique minimizer is $q_i^* = g$, the geometric-mean form~\eqref{eq:geometric_mean_target_app} below, and the minimum value is $F_i = -(1+B)\log Z_g \geq 0$, nonnegative because $Z_g \leq 1$ by H\"older's inequality (the exponents sum to one). The exponents $1/(1+B)$ are the reciprocal of the entropy weight discussed in Appendix~\ref{app:covariance_dynamics}: the best response is a Gibbs state at temperature $1+B$ relative to Bayesian updating. In a normalized joint the analogous coordinate update is a product of the prior and pair factors with unit exponents, and its stationarity includes each agent's appearance in its neighbors' factors~\cite[Ch.~5, local Gibbs bound and exact block E-coordinate]{DennisGaugeVFE2026}.

\subsection*{Conditional representation theorem}

\begin{theorem}[Conditional representation for the frozen-neighbor best response]\label{thm:representation_app}
Let\/ $\mathcal{D}(q,r) = \int r(c)f(q(c)/r(c))dc$ with $f$ \emph{real-analytic on $(0, \infty)$}, convex, $f(1)=0$, and, without loss of generality\footnote{Any $f$-divergence is invariant under $f(t)\mapsto f(t)+c(t-1)$ when $q$ and $r$ are normalized, so we may always choose the representative with $f'(1)=0$.}, $f'(1)=0$. Fix agent $i$'s neighbor beliefs $\{q_j\}$ and weights $\beta_{ij} \geq 0$ with $B := \sum_j \beta_{ij} > 0$, and let $\mathcal{D}$ enter the best-response problem~\eqref{eq:F_i_app} linearly. Then the stationary best response has the geometric-mean Boltzmann form
\begin{equation}
q_i^*(c) \propto p_i(c)^{1/(1+B)}\prod_{j}\bigl[\Omega_{ij}(c)q_j(c)\bigr]^{\beta_{ij}/(1+B)}
\label{eq:geometric_mean_target_app}
\end{equation}
for all choices of prior $p_i$ and neighbor beliefs $\{q_j\}$ if, and only if, $\mathcal{D}$ is the relative entropy (up to the affine equivalence above):
\[
\mathcal{D}(q,r) = D_{\mathrm{KL}}(q\|r).
\]
For $B = 1$, as when $\beta_i$ is a probability vector, the exponents are $1/2$.
\end{theorem}

\begin{corollary}[Closure forces the KL generator among nondegenerate divergences]
\label{cor:closure_forces_kl}
Within the same class, with real-analytic generator and nondegeneracy $f''(1) > 0$, exponential-family closure of the best response for all Gaussian priors and neighbor beliefs already forces $f'(t) = \gamma \log t$ with $\gamma > 0$, that is, $\mathcal{D} = \gamma D_{\mathrm{KL}}$, where the positive multiple $\gamma$ is absorbable into the coupling weights $\beta_{ij}$; the geometric-mean exponents of Eq.~\eqref{eq:geometric_mean_target_app} normalize $\gamma = 1$. The nondegeneracy hypothesis cannot be dropped: $f \equiv 0$ is real-analytic and convex, satisfies the normalization, and is closed for every Gaussian prior and neighbor set, since stationarity then reduces to $q_i^* = p_i$; it yields $\gamma = 0$ and the identically zero divergence. Because $f'(t) = \gamma \log t$ gives $f''(1) = \gamma$, the hypothesis $f''(1) > 0$ is exactly the statement $\gamma > 0$ and excludes that degenerate solution alone.
\end{corollary}

\emph{Proof sketch.} Stationarity gives $\log(q_i/p_i) + 1 + \sum_j \beta_{ij} f'(r_{ij}) = \lambda$. If $q_i^*$ is required to be log-quadratic for all Gaussian $(p_i, \{q_j\})$, then $h(x) := f'(e^x)$ must map every attained quadratic in $c$ to a quadratic in $c$; restricting a nondegenerate attained quadratic to lines and using the $t \mapsto -t$ symmetry of the restriction forces $h$ to be affine on the range of the quadratic, and real-analyticity extends $h$ affinely to all of $\mathbb{R}$, so $f'(t) = \alpha + \gamma\log t$. The normalization $f'(1) = 0$ kills $\alpha$, convexity gives $\gamma \ge 0$, and nondegeneracy gives $\gamma > 0$. Finally $\gamma D_{\mathrm{KL}}$ with weights $\beta_{ij}$ is identical to $D_{\mathrm{KL}}$ with weights $\gamma\beta_{ij}$. \hfill $\square$

\paragraph{Proof of Theorem~\ref{thm:representation_app}.}

\textbf{Step 1. Stationarity for a general $f$-divergence.} With ratio $r(c)=q_i(c)/\bigl[\Omega_{ij}(c)q_j(c)\bigr]$, the first-slot functional derivative is $\delta\mathcal{D}/\delta q_i(c) = f'(r(c))$, so the stationarity condition $\delta F_i/\delta q_i(c)=\lambda$, with Lagrange multiplier $\lambda$ for normalization, is
\begin{equation}
\log\frac{q_i(c)}{p_i(c)} + 1
+ \sum_j \beta_{ij}f'\left(\frac{q_i(c)}{\Omega_{ij}(c)q_j(c)}\right) = \lambda.
\label{eq:general_stationarity_app}
\end{equation}

\textbf{Step 2. Sufficiency.} For $\mathcal{D}=D_{\mathrm{KL}}$ the generator is $f(t)=t\log t - t + 1$ with $f'(t)=\log t$, and Lemma~\ref{lem:kl_sum_identity} gives the minimizer~\eqref{eq:geometric_mean_target_app} directly.

\textbf{Step 3. Necessity.} Assume the best response has the form~\eqref{eq:geometric_mean_target_app} for all priors $p_i$ and neighbor beliefs $\{q_j\}$, with general $B$. Taking logarithms, $(1+B)\log q_i^* = \log p_i + \sum_j\beta_{ij}\log(\Omega_{ij}q_j) + C$, and splitting $(1+B)\log q_i^* = \log q_i^* + \sum_j\beta_{ij}\log q_i^*$ gives
\begin{equation}
\log\frac{q_i^*}{p_i} + \sum_j\beta_{ij}\log\frac{q_i^*}{\Omega_{ij}q_j} = C',
\label{eq:rearranged_target_app}
\end{equation}
with $C' := -C$. It remains to pass from Eqs.~\eqref{eq:general_stationarity_app} and~\eqref{eq:rearranged_target_app} to the per-term identity
\begin{equation}
f'\left(\frac{q_i}{\Omega_{ij}q_j}\right) = \log\frac{q_i}{\Omega_{ij}q_j} + k,
\label{eq:ode_app}
\end{equation}
with a single constant $k$, and from there to $f'(t) = \log t + k$ on all of $(0, \infty)$. The Lagrange multiplier and the normalizer vary from configuration to configuration, and the geometric-mean form constrains the minimizer, not the neighbor beliefs. Both points are handled by fixing one configuration and sweeping the base point $c$ within it, where the multiplier and the normalizer are constants.

\begin{lemma}[Per-term separation and richness]
\label{lem:richness_app}
Let $f$ be real-analytic on $(0, \infty)$ with $f'(1) = 0$, and suppose the stationary solution of~\eqref{eq:general_stationarity_app} has the geometric-mean form~\eqref{eq:geometric_mean_target_app} for all priors and neighbor beliefs. Then $f'(t) = \log t$ as a functional identity on $(0, \infty)$.
\end{lemma}

\begin{proof}
Subtracting the rearranged target identity~\eqref{eq:rearranged_target_app} from the stationarity condition~\eqref{eq:general_stationarity_app} gives, for each fixed configuration $(p_i, \{q_j\}, \beta_i)$,
\begin{equation}
\sum_j \beta_{ij}\left[f'\big(r_{ij}(c)\big) - \log r_{ij}(c)\right] = \text{const},
\label{eq:summed_stationarity_app}
\end{equation}
with $r_{ij}(c) = q_i^*(c)/\big(\Omega_{ij}(c)q_j(c)\big)$ and a constant that depends on the configuration but not on $c$. To separate the terms, take a single-neighbor configuration: one neighbor $j$ with coupling $\beta_{ij} = B > 0$ and all other couplings zero. Then~\eqref{eq:summed_stationarity_app} reads $f'(r(c)) - \log r(c) = k$ directly, with one constant $k$ for the whole configuration; the Lagrange multiplier and the normalizer $Z_i$ are single constants within this fixed configuration and are absorbed into $k$. Under the geometric-mean form the single-neighbor ratio is $r(c) = Z_i^{-1}\big(p_i(c)/(\Omega_{ij}(c)q_j(c))\big)^{1/(1+B)}$, and choosing $p_i$ and $q_j$ so that the pointwise ratio $p_i(c)/(\Omega_{ij}(c)q_j(c))$ varies with $c$ (two non-identical Gaussians suffice) makes $r(c)$ sweep an open subinterval of $(0, \infty)$. The function $g(t) := f'(t) - \log t$ is real-analytic on $(0, \infty)$ and constant on that open subinterval, hence constant on all of $(0, \infty)$ by analytic continuation; this is the only point at which real-analyticity is used, and without it the identification holds on the attainable subinterval only. Thus $f'(t) = \log t + k$ on $(0, \infty)$, which is the per-term identity~\eqref{eq:ode_app} as a functional identity, and $f'(1) = 0$ forces $k = 0$.
\end{proof}

Integrating $f'(t) = \log t$ gives $f(t) = t\log t - t + C_0$, and $f(1)=0$ gives $C_0 = 1$. Therefore $f(t) = t\log t - t + 1$, the generator of $D_{\mathrm{KL}}(q\|r)$. \hfill $\square$

\subsection*{Scope: the joint scalar functional}
Theorem~\ref{thm:representation_app} concerns the best response alone. In the canonical scalar functional of Section~\ref{sec:mixture_derivation} each belief $q_i$ also sits in the second slot of every row $\beta_{ki}\mathcal{D}(q_k, \Omega_{ki}q_i)$ in which agent $i$ is a sender. The second-slot functional derivative of $\int r f(q/r)dc$ with respect to $r$ is $f(t) - tf'(t)$ at $t = q/r$, which for the relative entropy is $1 - t$. Free-form exponential-family closure of the joint stationarity condition would require both slot derivatives to be affine in $\log t$ on an open interval, $f'(t) = \alpha + a\log t$ and $f(t) - tf'(t) = \alpha' + c'\log t$. Differentiating, $f''(t) = a/t$ and $-tf''(t) = c'/t$, so $-a = c'/t$ on that interval, which forces $a = c' = 0$; then $f'' \equiv 0$, and the normalization with analyticity gives $f \equiv 0$. No nonzero $f$-divergence therefore gives free-form exponential-family closure once an agent is also a sender. A two-agent witness makes this concrete. In one dimension with $\Omega_{12} = \Omega_{21} = 1$, $\beta_{12} = \beta_{21} = 1$, $p_1 = \mathcal{N}(0, 1)$ and $p_2 = \mathcal{N}(3, 1/4)$ (second argument the variance), the geometric-pool laws $q_1 = \mathcal{N}(2, 1/2)$ and $q_2 = \mathcal{N}(8/3, 1/3)$ are mutual best responses, yet the gradient of the joint scalar at that point has norm $2.59$, and the joint Gaussian minimizer is $q_1 \approx \mathcal{N}(1.88, 0.83)$, $q_2 \approx \mathcal{N}(2.53, 0.60)$. In a Gaussian implementation the family imposes closure, so the closure criterion cannot select the divergence of the joint functional; the geometric-mean form~\eqref{eq:geometric_mean_target_app} describes the frozen-neighbor best response only.

\subsection*{Dual cost identification}
By the envelope theorem applied to the $\beta_{ij}$ optimization, the marginal attention cost is
\begin{equation}
C_{ij} \equiv \frac{\partial F_i}{\partial\beta_{ij}}\bigg|_{q_i^*}
= \mathcal{D}(q_i^*,\Omega_{ij}q_j),
\end{equation}
which for the relative entropy identifies the cost of attention with the divergence between the agent's best response and the neighbor's transported law. The identity holds for every divergence that enters linearly in $\beta_{ij}$, so it is a consistency property and does not select the divergence.

\subsection*{Relation to probabilistic opinion pooling}
The best response~\eqref{eq:geometric_mean_target_app} carries exponents $\beta_{ij}/(1+B)$ together with the prior exponent $1/(1+B)$, which sum to one over the prior-and-neighbor set, so it is the normalized weighted geometric (logarithmic) opinion pool, whose Gaussian case gives the precision-weighted aggregate $\Sigma_*^{-1} = \sum_k \lambda_k \Sigma_k^{-1}$. Dietrich and List~\cite{dietrich2016probabilistic} characterize this geometric pool axiomatically as externally Bayesian and unanimity/zero-preserving; Bordley~\cite{bordley1982multiplicative} and Genest and Zidek~\cite{GenestZidek1986} supply the multiplicative, log-linear precursor, which coincides with the normalized geometric pool in the continuous Gaussian fiber used here. Theorem~\ref{thm:representation_app} re-derives this operator, for the frozen-neighbor best response, by an $f$-divergence route within its hypothesis class. The meta-agent barycenter of Section~\ref{sec:meta_agent_definition} is the Gaussian M-projection of the linear pool that Dietrich and List characterize from eventwise independence. The axiomatic properties belong to the pool, not to its projection, and a form that drops the between-mean dispersion approximates the projection only in the high-coherence regime.

\subsection*{Interpretation and summary}
Gauge covariance fixes what is compared, the law $q_i$ and the transported law $\Omega_{ij}q_j$ in one frame, while the theorem above constrains how it is compared, and only for the best response. On the information-geometric side, the relative entropy is the Bregman divergence of the negative-entropy potential; among $f$-divergences on normalized distributions it is the only Bregman divergence~\cite{Amari2009alpha}, and it is the canonical divergence of the dually flat $e/m$ structure, whence the Pythagorean relation~\cite{AmariNagaoka2000}. The Fisher metric, by contrast, is the local Hessian of every smooth $f$-divergence, up to the factor $f''(1)$.

In Theorem~\ref{thm:representation_app} the competing $f$-divergences (the reverse relative entropy, $\chi^2$, squared Hellinger, and the Amari $\alpha$-divergences) are excluded already on the attainable subinterval of ratios, without invoking analyticity; a merely smooth generator agreeing with the relative entropy on that interval is observationally equivalent within the model. For Gaussian beliefs, exponential-family closure together with $f''(1) > 0$ already forces the KL generator up to a positive multiple absorbable into the coupling weights (Corollary~\ref{cor:closure_forces_kl}), so the geometric-mean target serves to normalize that multiple. By the $\mathrm{GL}(K_q)$ invariance theorem (Theorem~\ref{thm:glk_invariance}), the divergence is invariant under the full general linear group $\mathrm{GL}(K_q, \mathbb{R})$, which contains $\mathrm{SO}(K_q)$ and, after the realification $\mathrm{GL}(N, \mathbb{C}) \hookrightarrow \mathrm{GL}(2N, \mathbb{R})$ with $K_q = 2N$, the images of $\mathrm{U}(N)$ and $\mathrm{SU}(N)$.

